\documentclass[11pt]{article}
\usepackage[a4paper,margin=2.2cm]{geometry}
\usepackage{amsmath}
\setlength{\parindent}{0pt}
\setlength{\parskip}{5pt}
\newcommand{\un}[1]{\,\mathrm{#1}}

\begin{document}

{\Large\bfseries the duel, the maths behind what does work}

\vspace{4pt}
numbers are the films own unless i say its a guess

\section*{1. the reactor charging the flywheels}

the orbital ship has a 14 MW reactor that fills the flywheels over a day

\begin{align*}
E = P t = 14\times10^6 \times 24 \times 3600 = 14\times10^6 \times 86{,}400 = 1.21\times10^{12}\un{J} = 1.2\un{TJ}
\end{align*}

matches the films 1.2 TJ exactly. the reactor makes 140 MW of heat for 14 MW of power so 126 MW has to go out the radiators. $\varepsilon$ is emissivity, $\sigma$ is the Stefan Boltzmann constant, fins radiate off both sides

\begin{align*}
q &= 2\varepsilon\sigma T^4\\
\text{at } 800\un{K}: \quad q &= 2 \times 0.9 \times 5.67\times10^{-8} \times 800^4 = 41{,}800\un{W/m^2}\\
A &= \frac{126\times10^6}{41{,}800} = 3010\un{m^2}\\
\text{at } 600\un{K}: \quad A &= \frac{126\times10^6}{13{,}200} = 9530\un{m^2}
\end{align*}

so 55 by 55 m up to about 100 by 100 m of fins, easy on a 1.7 km ship. only thing a bit high is 10\% conversion, real thermoelectrics do 5 to 8\%

\section*{2. the flywheels}

56 wheels at 3000 rpm. radius is my guess, 3 m

\begin{align*}
E_{\text{wheel}} &= 1.2\times10^{12}/56 = 21.4\un{GJ}\\
\omega &= 3000 \times 2\pi/60 = 314\un{rad/s}, \qquad v_{\text{rim}} = \omega r = 314 \times 3 = 942\un{m/s}\\
m &= \frac{2E}{v^2} = \frac{2 \times 21.4\times10^9}{942^2} = 48{,}300\un{kg}
\end{align*}

hoop stress in a spinning rim, carbon fibre at 1600 kg/m$^3$

\begin{align*}
\sigma = \rho v^2 = 1600 \times 942^2 = 1.42\un{GPa}
\end{align*}

carbon fibre composite holds 2 to 3 GPa round the hoop so thats a real safety margin. half spin one way half the other so the spin cancels, $I$ is moment of inertia, $L$ angular momentum

\begin{align*}
I &= m r^2 = 48{,}300 \times 9 = 434{,}700\un{kg\,m^2}\\
L_{\text{one}} &= I\omega = 434{,}700 \times 314 = 1.37\times10^8\\
L_{\text{total}} &= 28 L - 28 L = 0
\end{align*}

if they all spun the same way thats $56 \times 1.37\times10^8 = 7.7\times10^9$ and turning the ship at 0.01 rad/s would take $\tau = L\Omega = 7.7\times10^7\un{N\,m}$ fighting it the whole time. the film got that right

\section*{3. the orbital gun shot}

rod is tungsten, 0.3 m across and 18 m long, barrel about 1016 m

\begin{align*}
m &= \rho \pi r^2 L = 19300 \times \pi \times 0.15^2 \times 18 = 24{,}560\un{kg}\\
KE &= \tfrac12 m v^2 = 0.5 \times 24{,}560 \times 5000^2 = 307\un{GJ} = 73\un{t\ TNT}\\
\eta &= 307/600 = 51\%
\end{align*}

all match the film. the long barrel keeps it gentle

\begin{align*}
a &= \frac{v^2}{2L} = \frac{5000^2}{2 \times 1016} = 12{,}300\un{m/s^2} = 1250\,g\\
t &= \frac{2L}{v} = \frac{2032}{5000} = 0.41\un{s}\\
P &= \frac{0.6\times10^{12}}{0.41} = 1.5\un{TW}
\end{align*}

the sabot grips the rod down its whole length so the load goes in as shear

\begin{align*}
F = m a = 24{,}560 \times 12{,}300 = 3.0\times10^8\un{N}, \qquad
\tau = \frac{F}{\pi d L} = \frac{3.0\times10^8}{\pi \times 0.3 \times 18} = 17.8\un{MPa}
\end{align*}

tungsten wont even notice that, compare 26 GPa on the destroyers round

\section*{4. the destroyers round going up}

to coast up to 380 km in a vacuum it needs

\begin{align*}
\tfrac12 v^2 = \mu\left(\tfrac1R - \tfrac1r\right) = 3.986\times10^{14} \times 8.835\times10^{-9} = 3.52\times10^6
\quad\Rightarrow\quad v = 2650\un{m/s}
\end{align*}

it has 7800 so its got plenty. leaving the air at 6.9 km/s at 80 km it gets to 380 km at

\begin{align*}
v^2 = 6900^2 - 2\mu\left(\tfrac{1}{6451\un{km}} - \tfrac{1}{6751\un{km}}\right) = 4.76\times10^7 - 5.49\times10^6 \quad\Rightarrow\quad v = 6490\un{m/s}
\end{align*}

nose heating at sea level (Sutton Graves), $r_n$ is nose radius, guess 3 mm

\begin{align*}
q = k\sqrt{\rho/r_n}\,v^3 = 1.74\times10^{-4} \times \sqrt{1.225/0.003} \times 7800^3 = 1.7\times10^{9}\un{W/m^2}
\end{align*}

tungsten has to be heated to 3695 K and melted before it comes off

\begin{align*}
H = c\Delta T + L_f = 134 \times 3400 + 192{,}000 = 0.65\un{MJ/kg}
\end{align*}

air density falls off with height ($H_s = 8.5$ km) and it climbs at $7800\sin70^\circ = 7330$ m/s, so the air it goes through adds up to

\begin{align*}
\int\!\sqrt{\rho}\,dt = \sqrt{\rho_0}\,\frac{2H_s}{v_z} = 1.107 \times \frac{17{,}000}{7330} = 2.57\\
\text{nose lost} = \frac{k v^3}{\sqrt{r_n}} \times \frac{2.57}{\rho_w H} = \frac{1.51\times10^9 \times 2.57}{19300 \times 0.65\times10^6} = 0.31\un{m}
\end{align*}

thats the worst case, the nose blunts as it melts which cuts the heating, so 10 to 30 cm off a 2 m rod. it survives and the white hot sheath in the film is right. nothing burns once its out of the air

\section*{5. the rod coming down and going in}

$\beta$ is ballistic coefficient, $C_d$ about 0.15

\begin{align*}
\beta = \frac{m}{C_d A} = \frac{24{,}560}{0.15 \times \pi \times 0.15^2} = 2.3\times10^6\un{kg/m^2}
\end{align*}

thats so heavy for its size the air barely touches it. same heating sum as above with a blunter 5 cm nose at 5 km/s

\begin{align*}
\int\!\sqrt{\rho}\,dt = 1.107 \times \frac{17{,}000}{5000\sin44^\circ} = 1.107 \times \frac{17{,}000}{3473} = 5.42\\
\text{nose lost} = \frac{1.74\times10^{-4} \times 5000^3}{\sqrt{0.05}} \times \frac{5.42}{19300 \times 0.65\times10^6} = 0.04\un{m}
\end{align*}

4 cm off 18 m so it hits whole. then in the water a long rod wears away as it goes, $L$ is rod length

\begin{align*}
P = L\sqrt{\rho_{\text{rod}}/\rho_{\text{water}}} = 18 \times \sqrt{19.3/1.025} = 18 \times 4.34 = 78\un{m}
\end{align*}

so it really does dive about 80 m and eat itself to nothing like the film shows, and dumping 73 t of TNT worth of energy right under the keel breaks the destroyer

\end{document}
